\section{Electrostatic stress and the asymptotically flat attachment}
\label{sec:stress-extension}

Throughout this section, $(\Omega,g)$ is a smooth connected exterior with
the topology, completeness, spherical end, tensor decay, and weighted
scalar integrability specified in Section~\ref{sec:hypotheses}.  We assume
\begin{equation}
 R_g\ge-6,\qquad H_g(S)\le0,
 \label{eq:stress-base-assumptions}
\end{equation}
and make no causal assumption on its mass covector.  The normal on $S$
points into $\Omega$.  Fix the end chart, and put
\begin{equation}
 A_* = A_{\min}(S;g),\qquad a=\sqrt{A_*/(4\pi)}.
 \label{eq:stress-area-radius}
\end{equation}
The objective is to construct a compact inner region and a scalar-flat
\AF\ exterior whose mass is at most $p_0(g)-a^3/2$, up to an arbitrarily
small error.  The tensor introduced for this purpose is synthetic: it is
not the second fundamental form of the original data.  We denote the
electrostatic stress by $\Pi$, reserving $\pi$ for the graph second form in
the subsequent construction.

\subsection{A conformal approximation of the end}

\begin{proposition}[End approximation]
\label{prop:conformal-approximation}
Under \eqref{eq:stress-base-assumptions}, there are smooth metrics $g_j$
on the same exterior such that $R_{g_j}>-6$, $H_{g_j}(S)\le0$,
$g_j\to g$ locally in $C^2$, and
\begin{equation}
 (1-\epsilon_j)g\le g_j\le(1+\epsilon_j)g,
 \quad \epsilon_j\downarrow0,
 \qquad p_0(g_j)\longrightarrow p_0(g).
 \label{eq:stress-density-convergence}
\end{equation}
Each $g_j$ has the required weighted scalar integrability and finite mass
fluxes.  Outside a compact set, depending on $j$, it is of the form
\begin{equation}
 g_j=\phi_j^4 b,\qquad
 \phi_j=1+\alpha_j(\omega)r^{-3}+O(r^{-3-\xi}),
 \label{eq:stress-conformal-end}
\end{equation}
where $\xi>0$ can be fixed and the remainder has the same decay after
any fixed finite number of angular and radial-logarithmic derivatives.
In particular $A_{\min}(S;g_j)\to A_{\min}(S;g)$.
\end{proposition}

\begin{proof}
Choose $j_0$ with $3/2<j_0<\tau$, where $\tau$ is the original end
decay exponent.  Extend $r$ to a smooth positive function on the compact
part.  For large $B$, let $\chi_B$ be a smooth radial cutoff equal to one
for $r\le B$ and zero for $r\ge2B$, with a fixed profile in $r/B$.
Set
\[
 g_B=b+\chi_B(g-b)
\]
on the end and $g_B=g$ on the compact part.  For sufficiently large $B$
this is positive definite, and its uniform metric charts and
$C^{2,\alpha}$ bounds can be chosen independently of $B$.  Fix a smooth
$k>0$ equal to $r^{-4}$ far out.  We seek a small $u$ with
$\partial_\nu u=0$ on $S$ satisfying
\begin{equation}
 -8\Delta_{g_B}u+(R_{g_B}+6)(1+u)
 +6\big((1+u)^5-(1+u)\big)
 =\chi_B(R_g+6)+\delta k.
 \label{eq:stress-density-equation}
\end{equation}
The constant-term discrepancy, excluding $\delta k$, is supported in
$B<r<2B$ and is $O(r^{-\tau})$ in local $C^{0,\alpha}$ norms.
Consequently it tends to zero in the scalar weighted norm of weight
$r^{-j_0}$.  Also $R_{g_B}+6\ge-o(1)$ uniformly, because it equals the
nonnegative $R_g+6$ on the inner part, vanishes beyond $2B$, and is
$O(B^{-\tau})$ on the intervening annulus.

For completeness, the linearized operator divided by eight is
\[
 \mathscr L_B=-\Delta_{g_B}+3+(R_{g_B}+6)/8.
\]
It has an inverse, with homogeneous Neumann condition, bounded uniformly
from weighted $C^{0,\alpha}$ to weighted $C^{2,\alpha}$ at weight
$j_0$ (a smaller H\"older exponent may be used).  To prove the required
supremum bound, take a positive function $\beta$ constant on a large
compact set, comparable to $r^{-j_0}$ at infinity, and let its logarithmic
slope decrease slowly from zero to $-j_0$ as a function of $\log r$.
For the limiting radial operator,
\[
 (-\partial_\eta^2-2\partial_\eta+3)e^{-j_0\eta}
 =(3+2j_0-j_0^2)e^{-j_0\eta}>0.
\]
Every intermediate slope has the same strict sign.  Slow variation of
the slope, placement of the transition sufficiently far out, and the
uniformly small negative part of $(R_{g_B}+6)/8$ therefore give
$\mathscr L_B\beta\ge c\beta$ with $c>0$ independent of large $B$.
Here $\partial_\nu\beta=0$ on $S$.  Solve on compact truncations with
zero outer Dirichlet data.  The maximum principle, applied to positive
and negative multiples of $\beta$, bounds the solution by the weighted
norm of its forcing.  Uniform local interior and Neumann Schauder
estimates give the stated second-order bound.  Compactness gives a
solution on $\Omega$, and the maximum principle gives uniqueness.
This proves the inverse assertion without an assumption about solvability
on the noncompact manifold.

The nonlinear remainder in \eqref{eq:stress-density-equation} is
$O(u^2)$ and is locally Lipschitz with Lipschitz constant tending to zero
on small weighted balls.  The uniform inverse and contraction mapping
therefore give $u=u_{B,\delta}$ tending to zero in weighted
$C^{2,\alpha}_{j_0}$ as $B\to\infty$ and $\delta\downarrow0$.
Local elliptic bootstrapping makes $u$ smooth.  For small parameters
$1+u>0$, and the conformal scalar-curvature formula gives, for
$\widetilde g=(1+u)^4g_B$,
\begin{equation}
 R_{\widetilde g}+6
 =(1+u)^{-5}\big(\chi_B(R_g+6)+\delta k\big)>0.
 \label{eq:stress-density-scalar}
\end{equation}
The boundary conformal formula and $\partial_\nu u=0$ give
$H_{\widetilde g}=(1+u)^{-2}H_g\le0$.  The uniform metric convergence in
\eqref{eq:stress-density-convergence} follows from the weighted bound and
the cutoff construction.

We next justify the end expansion for each fixed $B,\delta$.
Beyond $2B$, $g_B=b$ and
\begin{equation}
 (-\Delta_b+3)u=F(u)+\delta r^{-4}/8,
 \qquad F(u)=O(u^2),\quad F'(u)=O(u).
 \label{eq:stress-density-end-ode}
\end{equation}
All fixed orders of angular derivatives of $u$ are $O(r^{-j_0})$.
Indeed rotations commute with $\Delta_b$.  Their difference quotients
solve a linear equation with small potential $O(u)$, with smooth data
on a fixed inner sphere and decaying data at infinity.  The same
$r^{-j_0}$ barriers and local weighted estimates bound these quotients.
Induction, with the lower angular derivatives as forcing, proves the
claim at every fixed order.  Radial derivatives at the initial weighted
order follow from local elliptic estimates.  Since
\[
 \Delta_b=\partial_\eta^2+2\coth\eta\,\partial_\eta
            +r^{-2}\Delta_\sigma,
\]
Equation~\eqref{eq:stress-density-end-ode} reduces, with all these angular
derivatives, to
\[
 -u_{\eta\eta}-2u_\eta+3u=O(e^{-(3+\xi)\eta})
\]
for any fixed sufficiently small
$0<\xi<\min\{2j_0-3,j_0-1,1\}$.  Variation of constants for the two
solutions $e^\eta,e^{-3\eta}$ excludes the growing mode by the known
decay and gives $u=c(\omega)e^{-3\eta}+O(e^{-(3+\xi)\eta})$.
The same argument after angular differentiation makes $c$ smooth;
the equation gives all stated radial derivatives.  Changing from
$\eta$ to $r$ gives \eqref{eq:stress-conformal-end}.

We finally check mass convergence, since convergence only in a weight
less than three would not suffice.  Write $h=b+e$ for either the original
metric or one of the approximants, and let $\mathbb U_{V_0}(e)$ be the
one-form in the mass definition.  The linearized scalar curvature is
\[
 (DR)_b(e)=\divg_b\divg_b e-\Delta_b\tr_b e+2\tr_b e.
\]
Differentiation, using $\nabla_b^2V_0=V_0b$, yields the exact linear
identity $\divg_b\mathbb U_{V_0}(e)=V_0(DR)_b(e)$.
The difference between $R_h+6$ and $(DR)_b(e)$ is bounded by
$C(|e|\,|\nabla_b^2e|+|\nabla_be|^2+|e|^2)$.  The uniform weighted
second-order bound and $dV_b\asymp r\,dr\,dA_\sigma$ thus give
\begin{equation}
 16\pi p_0(h)-\int_{r=D}\mathbb U_{V_0}(h-b)(\nu_b)\,\dd A_b
 =\int_{r>D}V_0(R_h+6)\,\dd V_b+O(D^{3-2j_0}),
 \label{eq:stress-density-mass-tail}
\end{equation}
uniformly in the approximants.  The scalar integrands in
\eqref{eq:stress-density-scalar} are dominated by a fixed multiple of
$V_0(R_g+6+k)$, which is integrable; they converge pointwise to
$V_0(R_g+6)$.  The same divergence calculation with $V_i$, using
$|V_i|\le V_0$ and the corresponding derivative bounds, proves existence
of all fluxes for the approximants.  At a fixed $D$ their surface fluxes
converge by local $C^2$ convergence.  Dominated convergence in
\eqref{eq:stress-density-mass-tail}, followed by $D\to\infty$, proves
the asserted time-component convergence.  This is also the scalar
linearization underlying the mass definition in \citet{CH2003}.
Finally, the two-sided metric comparison compares the areas of every
admissible intrinsic boundary with the same multiplicative factors;
taking infima proves convergence of the minimum enclosing areas.
\end{proof}

We henceforth fix one approximating metric, write it again as $g$, and
write $\alpha$ for its smooth coefficient in
\eqref{eq:stress-conformal-end}.  Thus $R_g+6>0$ everywhere.  All
subsequent constants may depend on this fixed metric; no uniformity
across the approximation is needed.

\subsection{The nonlinear electrostatic problem}

For $0\le e\le1$ define
\begin{equation}
 t(e)=\sqrt{1-e^{3/2}},\qquad
 L(w)=\max_{0\le e\le1}\{t(e)+we\},\qquad w\ge0.
 \label{eq:stress-legendre-law}
\end{equation}
Strict concavity of $t$, together with $t'(0)=0$ and
$t'(e)\to-\infty$ as $e\uparrow1$, gives a unique maximizing $e=e(w)$,
with $e(0)=0$, $0<e(w)<1$ for $w>0$, and
\begin{equation}
 w=-t'(e)=\frac{3\sqrt e}{4t(e)},\qquad
 L'(w)=e(w),\qquad L=t+we.
 \label{eq:stress-constitutive-law}
\end{equation}
In particular $L(|\xi|)$ is convex as a function of a covector $\xi$.

\begin{proposition}[Electrostatic potential, stress, and charge]
\label{prop:electrostatic-stress}
For every $M_e>0$ there is a weak solution $u$ of
\begin{equation}
 \divg_g E=0,\qquad
 E=e(w)\frac{\nabla u}{w},\quad w=|du|_g,
 \qquad u|_S=-M_e,\quad u\longrightarrow0\text{ at infinity},
 \label{eq:stress-electrostatics}
\end{equation}
where $E=0$ at $w=0$.  It satisfies $-M_e\le u\le0$,
$|u|+|du|_g=O(r^{-1})$, and is $C^{1,\beta}$ locally up to $S$ and
in compactified coordinates $(\rho,\omega)$ up to $\rho=0$, for some
$\beta>0$, where $\rho=e^{-\eta}$.
The continuous symmetric tensor
\begin{equation}
 \Pi=L(w)g-du\otimes E^\flat
 \label{eq:stress-tensor}
\end{equation}
is distributionally divergence free, is positive definite, and obeys
\begin{equation}
 (\tr_g\Pi)^2-2|\Pi|_g^2=3.
 \label{eq:stress-algebraic-identity}
\end{equation}
There is a continuous nonnegative function $e_0$ on $\Sph^2$ such that,
uniformly in angle,
\begin{equation}
 r^2|E|_g\longrightarrow e_0,\qquad
 r^3\big(\Pi(\nu,\nu)-1\big)
       \longrightarrow-\tfrac12 e_0^{3/2}.
 \label{eq:stress-end-asymptotics}
\end{equation}
Here $\nu$ points toward increasing $r$.  The charge
\begin{equation}
 F_{M_e}=\int_{r=R}g(E,\nu)\,\dd A_g
 \label{eq:stress-charge-definition}
\end{equation}
is independent of large $R$, satisfies
$F_{M_e}\le\int_{\Sph^2}e_0\,\dd A_\sigma$, and
\begin{equation}
 \liminf_{M_e\to\infty}F_{M_e}\ge A_*.
 \label{eq:stress-charge-saturation}
\end{equation}
\end{proposition}

\begin{proof}
We give the construction and its regularity before establishing the
stress and charge assertions.

\emph{Constitutive estimates and regularization.}
For $w>0$ put $Q(w)=we'(w)/e(w)$.  Logarithmic differentiation of
\eqref{eq:stress-constitutive-law} gives
\begin{equation}
 Q(w)^{-1}=\frac12+\frac{3e(w)^{3/2}}{4(1-e(w)^{3/2})}.
 \label{eq:stress-Q-formula}
\end{equation}
Thus $0<Q\le2$.  At zero, inversion in the variable $\sqrt e$ gives
\[
 e(w)=\frac{16}{9}w^2+O(w^5);
\]
in particular $e(w)/w^2$ extends smoothly and positively to $w=0$.
At infinity, inversion in $w^{-2}$ gives
\begin{equation}
 e(w)=1-\frac{3}{8w^2}+O(w^{-4}),\qquad
 Q(w)=\frac{3}{4w^2}+O(w^{-4}),\qquad
 wQ'(w)=-\frac{3}{2w^2}+O(w^{-4}).
 \label{eq:stress-large-gradient-law}
\end{equation}
Regularize only at zero by
\begin{equation}
 e_\varepsilon(w)=w\frac{e(z)}z,\quad
 z=(w^2+\varepsilon^2)^{1/2},\quad
 Q_\varepsilon(w)=1+\frac{w^2}{z^2}(Q(z)-1),\quad 0<\varepsilon<1.
 \label{eq:stress-regularized-law}
\end{equation}
The flux $e_\varepsilon(|\xi|)\xi^\sharp/|\xi|$ is smooth and strictly
elliptic on every finite gradient range, including zero.  Its radial
and transverse derivative eigenvalues are $e_\varepsilon'$ and
$e_\varepsilon/w$, respectively.  We have $0<Q_\varepsilon\le2$;
uniformly for sufficiently large $w$,
\begin{equation}
 c\le w^2Q_\varepsilon(w)\le C,\qquad
 -w^2\big(wQ_\varepsilon'(w)+Q_\varepsilon(w)\big)\ge c.
 \label{eq:stress-uniform-Q}
\end{equation}
These follow by substituting \eqref{eq:stress-large-gradient-law} in
\eqref{eq:stress-regularized-law}; the leading coefficient is
$3/4+\varepsilon^2$, bounded above and below independently of
$\varepsilon$.

\emph{Boundary and interior gradient bounds.}
On a connected smooth truncation $\Omega_R$ prescribe $u=-M_e$ on $S$
and $u=0$ on $r=R$.  The maximum principle gives $-M_e\le u\le0$.
At points with nonzero gradient the regularized equation is
\begin{equation}
 P:\nabla^2u=0,\qquad
 P=g^{-1}+(Q_\varepsilon(w)-1)n\otimes n,
 \quad n=\nabla u/w.
 \label{eq:stress-regularized-nondivergence}
\end{equation}
Let $h$ be inward distance from $S$.  Since $H(S)\le0$, smoothness
gives $\Delta h\le C h$ in a fixed collar.  A supersolution above the
inner boundary datum is $-M_e+B(h)$, where
\[
 B(h)=b_0^{-1}\log(1+Dh),\qquad B''=-b_0(B')^2.
\]
For $p=B'$ sufficiently large, \eqref{eq:stress-uniform-Q} gives
\[
 P:\nabla^2B=Q_\varepsilon(p)B''+p\Delta h
 \le-cb_0+C/b_0,
\]
because $hp\le1/b_0$.  Choose $b_0$ first to make this negative,
then a collar width $h_0$ so small that
$(2b_0h_0)^{-1}$ is above the required gradient threshold, and finally
$D$ large so that $B(h_0)>M_e$ and $B'\ge(2b_0h_0)^{-1}$ there.
Comparison with the constant lower barrier bounds the inward derivative
on $S$, independently of $R$ and $\varepsilon$.
At the outer sphere, inward distance may be replaced by
$h=\eta_R-\eta$ in a uniform collar.  For gradients parallel to
$d\eta$, asymptotic hyperbolicity and $Q_\varepsilon\le2$ give
$P:\nabla^2h\le-1$ for sufficiently large $R$.  The linear lower
barrier $-C h$, together with $u\le0$, bounds the outer derivative
with $C$ depending on $M_e$ and the fixed collar width only.

For the interior estimate use an orthonormal frame with first vector
$n$, and write transverse indices as $a,b$.  Differentiating
\eqref{eq:stress-regularized-nondivergence}, commuting one derivative,
and differentiating the norm of $du$ gives
\begin{align}
 P:\nabla^2\log w
 ={}&\Ric(n,n)
 +(1-Q_\varepsilon)\frac{|d_\perp w|^2}{w^2}
 +\frac{\sum_{a,b}u_{ab}^2}{w^2} \notag\\
 &-\big(wQ_\varepsilon'+Q_\varepsilon\big)\frac{w_1^2}{w^2}.
 \label{eq:stress-gradient-identity}
\end{align}
To specify the cancellation, the differentiated coefficient term is
$Q_\varepsilon'w_1^2+2(Q_\varepsilon-1)|d_\perp w|^2/w$;
the differentiated norm contributes
$[Q_\varepsilon\sum_a u_{1a}^2+\sum_{a,b}u_{ab}^2]/w$ before
division by $w$.  Subtracting $P(dw,dw)/w^2$ produces exactly
\eqref{eq:stress-gradient-identity}.
At a positive interior maximum of $w\exp(B_0u)$ we have
$d_\perp w=0$ and $w_1/w=-B_0w$.  The left side of
\eqref{eq:stress-gradient-identity} is nonpositive, since
$P:\nabla^2u=0$.  The global lower Ricci bound and
\eqref{eq:stress-uniform-Q} make the right side at least
$-C+cB_0^2$ when $w$ is large.  A fixed sufficiently large $B_0$
excludes such a maximum.  Combining the remaining interior maxima
with the boundary estimates and the height bound gives
\begin{equation}
 \sup_{\Omega_R}|du|_g\le C(M_e),
 \label{eq:stress-global-gradient-bound}
\end{equation}
uniformly in large $R$ and $0<\varepsilon<1$.

\emph{Existence on truncations and passage to the limit.}
For fixed $R,\varepsilon$, vary the inner Dirichlet value from zero
to $-M_e$.  The derivative of the divergence equation is a strictly
elliptic divergence operator without zeroth-order term and is
invertible with zero Dirichlet data by the maximum principle and
linear elliptic theory.  The bounds already proved hold on this
entire path.  On its bounded gradient range, ordinary uniformly
elliptic quasilinear gradient H\"older estimates apply.  For the
boundary estimate subtract its constant value and reflect oddly in
boundary normal coordinates, reflecting the metric as well.  The
reflected metric is Lipschitz; its mixed normal--tangential entries
vanish on the face.  The reflected equation is weakly valid by the
zero trace.  Equivalently, difference quotients give uniformly
elliptic equations for the first derivatives with bounded divergence
forcing.  Scalar H\"older estimates followed by Schauder estimates
on the original smooth side give closedness of the continuation.
Here constants may depend on the fixed $R,\varepsilon$;
\citet{GT2001} supplies the ordinary elliptic estimates being used.
Thus smooth regularized solutions exist.

A decay bound is uniform in both regularizations.  On the end put
$v=C(e^{-\eta}-e^{-2\eta})$.  It is positive for large $\eta$, and
\[
 \frac{v''}{-v'}=\frac{1-4e^{-\eta}}{1-2e^{-\eta}}
 \le1-e^{-\eta}.
\]
Also $|d\eta|_g^2=1+O(e^{-\tau\eta})$ and
$P:\nabla^2\eta=2+O(e^{-\min\{\tau,2\}\eta})$ for the
radial test direction, uniformly for $0<Q_\varepsilon\le2$.
Since $\tau>1$, these inequalities imply
$P:\nabla^2v\le0$ sufficiently far out, independently of $C$.
Comparison of the nonnegative solution $-u$ with $v$, choosing $C$
on a fixed inner sphere and using $u=0$ on the truncating edge,
gives $|u|\le Ce^{-\eta}$.

Take a subsequence with $R\to\infty$ and $\varepsilon\downarrow0$.
The uniform Lipschitz bound gives local uniform convergence and weak
gradient convergence to a locally Lipschitz function with the required
traces.  To pass to the equation, let $\chi$ be a compactly supported
cutoff and test the regularized equation with $\chi(u_j-u)$.
The terms containing $d\chi$ tend to zero by uniform convergence.
Subtracting the limiting-law flux evaluated at $du$ gives
\[
 \int\chi\,
 \big(A(du_j)-A(du)\big)\cdot(du_j-du)\,\dd V_g\longrightarrow0,
 \qquad A(\xi)=e(|\xi|)\xi^\sharp/|\xi|.
\]
Uniform convergence of the constitutive functions on bounded ranges
justifies replacing the regularized flux here.  The law $A$ is
strictly monotone.  On the fixed compact gradient range its
monotonicity pairing has a positive minimum whenever the two
arguments are separated by a prescribed positive distance.
Thus $du_j\to du$ in measure and strongly in every finite local
$L^p$.  This proves the limiting weak equation.

\emph{Gradient regularity and compactification.}
We apply the scalar $p=3$ gradient theorem recorded in
Section~\ref{sec:preliminaries}, namely Theorem~1.1 and estimate~(1.8)
of \citet{AraujoZhang2020}.  The coordinate flux must satisfy globally
\begin{align}
 \lambda|\xi|\,|\zeta|^2
 &\le \mathcal A_\xi(x,\xi)[\zeta,\zeta],
 &|\mathcal A|+|\mathcal A_\xi|\,|\xi|&\le C|\xi|^2,
 \label{eq:stress-p3-structure}\\
 |\mathcal A(x,\xi)-\mathcal A(y,\xi)|
 &\le C|x-y|\,|\xi|^2.
 \notag
\end{align}
It must also be $C^1$ in $\xi$ and the solution must have bounded
local $W^{1,3}$ energy.  The Lipschitz spatial bound supplies a
$C^{0,1/2}$ modulus, so $n=p=3$ and zero forcing meet all the
parameters of that theorem.  Throughout this argument choose only
$0<\beta<\min\{\alpha_m,1/4\}$ as in
Section~\ref{sec:preliminaries}; the decay and boundary limits below
require a positive exponent, not an optimal one.

We verify global, rather than merely bounded-gradient, structure before
applying this theorem.  On the gradient range in
\eqref{eq:stress-global-gradient-bound}, $e(w)/w^2$ and $Q(w)$ have
positive upper and lower bounds.  Thus $e'(w)$ is comparable to $w$.
Choose a smooth cutoff $\chi=1$ below the known range and zero above
twice a larger fixed bound, and replace $e'$ outside that range by
\[
 \widetilde e'(w)=\chi(w)e'(w)+(1-\chi(w))c_1w,
 \qquad \widetilde e(0)=0,
\]
where $c_1>0$.  Then $\widetilde e=e$ on the range attained by the
solution, $\widetilde e'\asymp w$, and $\widetilde e\asymp w^2$
globally.  In coordinates with metric matrix $G$ the density flux is
\begin{equation}
 \mathcal A^i(x,\xi)=\sqrt{\det G}\,
 \frac{\widetilde e(|\xi|_G)}{|\xi|_G}\,G^{ij}\xi_j,
 \qquad |\xi|_G=(G^{ij}\xi_i\xi_j)^{1/2}.
 \label{eq:stress-coordinate-flux}
\end{equation}
Uniform positive metric bounds give
\eqref{eq:stress-p3-structure}; Lipschitz metric coefficients give
the spatial bound.  The flux is $C^1$ at zero because it is
$O(|\xi|^2)$ and its derivative is $O(|\xi|)$.  Its derivative is
symmetric, as is also apparent from its convex radial primitive.
The solution has finite local cubic energy by its Lipschitz bound.
The quoted theorem therefore proves local $C^{1,\beta}$ regularity.

The same argument works at $S$.  In orthogonal boundary coordinates
with reflection matrix $J=\operatorname{diag}(-1,1,1)$, extend the
metric matrix by $G(-h,x)=JG(h,x)J$, and extend $u+M_e$ oddly.
The mixed metric entries vanish on the face, so this extension is
Lipschitz.  The reflected flux satisfies
$\mathcal A_-(\xi)=J\mathcal A_+(J\xi)$.  Evaluated on the odd
solution, its normal component is even and its tangential components
are odd.  A test on the reflected ball consequently reduces to a
test on the original half-ball with zero trace.  Such tests are
allowed: cutoff at distance $h=\delta$ creates an error tending to
zero since the test is $O(h)$, the flux is bounded, and the strip
volume tends to zero.  Thus the reflected weak equation satisfies
all the same hypotheses, proving the asserted boundary regularity.

To obtain gradient decay, take a uniform metric ball far out with
central radius $r_*$, and put $s_*=C/r_*$, so that $|u|\le s_*$ on a
slightly larger ball.  Apply the preceding fixed global extension
to the equation for $u/s_*$; its flux is
\[
 \mathcal A_{s_*}(x,\xi)=s_*^{-2}\mathcal A(x,s_*\xi).
\]
The cubic structure constants and spatial modulus are unchanged.
The required energy bound does not presume a bound on the normalized
gradient.  Testing with $\chi^3(u/s_*)$, and using coercivity and
Young's inequality, gives
\begin{equation}
 \int\chi^3|d(u/s_*)|^3
 \le C\int |u/s_*|^3|d\chi|^3\le C.
 \label{eq:stress-scaled-energy}
\end{equation}
The H\"older estimate for the gradient, together with this energy
bound, gives its supremum on a smaller ball: some point has gradient
bounded by the cubic mean, and the H\"older bound controls its
difference from every other point.  Therefore $|du|_g\le C/r$.

Now let $g_c=\rho^2g$ on the end.  It extends to a positive Lipschitz
metric up to $\rho=0$, with vanishing mixed normal--tangential entries
there.  In the conformal end at hand this follows directly from
\[
 g_c=\phi^4\left(d\rho^2+\tfrac14(1-\rho^2)^2\sigma\right)
\]
and \eqref{eq:stress-conformal-end}.  The estimate just obtained gives
$|du|_{g_c}\le C$, and the height estimate gives the zero trace.
Writing $\bar w=|du|_{g_c}$, the coordinate divergence equation in
this metric has radial law
\begin{equation}
 q_\rho(\bar w)=\rho^{-2}e(\rho\bar w)
              =\bar w^2 F(\rho\bar w),
 \qquad F(t)=e(t)/t^2.
 \label{eq:stress-compactified-law}
\end{equation}
On the known bounded range of $\bar w$, this satisfies uniform cubic
structure and Lipschitz spatial dependence, including at $\rho=0$.
For an explicit uniform extension, set
$H(t)=2F(t)+tF'(t)>0$, choose a fixed cutoff equal to one below that
range and zero above twice a larger bound, and define
\[
 \widetilde q_\rho'(w)
 =w\big[\chi(w)H(\rho w)+(1-\chi(w))c_1\big],
 \qquad \widetilde q_\rho(0)=0.
\]
It agrees with $q_\rho$ on the attained range, is comparable to
$w^2$, has derivative comparable to $w$, and has the required uniform
Lipschitz dependence in position.  Reflect $g_c$ as above, use
$|\rho|$ in this coefficient, and reflect $u$ oddly.  The gradient
bound ensures $W^{1,3}$ across the face.  The same zero-trace test
argument proves the reflected weak equation.  The scalar theorem
now gives $C^{1,\beta}$ in compactified coordinates.  No theorem
about degenerate boundary equations is being assumed in addition
to the stated interior result.

\emph{Stress and asymptotic coefficients.}
Convexity shows that the weak solution minimizes the integral of
$L(|du|)$ against compactly supported Lipschitz variations: integrate
the convex tangent inequality and use \eqref{eq:stress-electrostatics}.
Apply this fact to compactly supported changes of independent variable.
For the flow of a smooth vector field $X$, differentiating the local
energy of the transported potential gives
\[
 \int\big(L(w)g-du\otimes E^\flat\big):\nabla X\,\dd V_g=0.
\]
The calculation is justified by bounded gradients on the support;
no finite total energy of $L$ is needed.  This is
$\divg_g\Pi=0$ in distributions.
The stress eigenvalues in the gradient and transverse directions are
$t$ and $l=L=t+we$, respectively, with $\Pi=g$ at zero gradient.
They are positive at every finite gradient.  Moreover,
\[
 (t+2l)^2-2(t^2+2l^2)=4tl-t^2
 =3t^2+4twe=3(1-e^{3/2})+3e^{3/2}=3,
\]
proving \eqref{eq:stress-algebraic-identity}.

Since $u$ has constant trace at conformal infinity, all its tangential
derivatives in compactified coordinates vanish there.  Its scaled
gradient consequently has a continuous radial limit.  Together with
$r\rho\to1/2$ and $e(w)=(16/9)w^2+O(w^5)$, this gives a continuous
$e_0\ge0$ with $r^2e(w)\to e_0$ uniformly.  At a point where
$e_0>0$ the gradient direction approaches the radial direction, so
$\Pi(\nu,\nu)=t+o(r^{-3})$.  At a point where $e_0=0$, both
$t-1$ and $l-t=we$ are $O(e^{3/2})=o(r^{-3})$, and the same conclusion
holds.  Uniformity follows by splitting into the sets where $e_0$
is larger or smaller than an arbitrary fixed positive threshold.
Since $t=1-e^{3/2}/2+O(e^3)$, this proves
\eqref{eq:stress-end-asymptotics}, as well as
$E=O(r^{-2})$ and $\Pi-g=O(r^{-3})$ in metric norm.

\emph{Flux saturation.}
Continuous distributionally divergence-free fields have the usual
flux identities on smooth compact domains: use smooth cutoffs
approaching their boundaries and continuity in the shrinking collars.
This proves independence in \eqref{eq:stress-charge-definition}.
The asymptotics and $dA_g/r^2\to dA_\sigma$ give
$F_{M_e}\le\int e_0\,dA_\sigma$.  Integration by parts in the
electrostatic equation gives
\begin{equation}
 M_eF_{M_e}=\int_\Omega w e(w)\,\dd V_g.
 \label{eq:stress-energy-charge}
\end{equation}
Indeed the infinity term tends to zero because
$u=O(r^{-1})$, $E=O(r^{-2})$, and the sphere area is $O(r^2)$.
The inner term has the indicated sign: the outward normal of the
exterior at $S$ is $-\nu$, while $u=-M_e$ there.  The integral is
finite since $we=O(r^{-3})$ and $dV_g\asymp r\,dr\,dA_\sigma$.

Fix a large $R_0$, choose $4/3<s_0<2$, and set
\begin{equation}
 \zeta=1\quad(r\le R_0),\qquad
 \zeta=(R_0/r)^{s_0}\quad(r>R_0),\qquad
 C_\zeta=\int_\Omega(1-t(\zeta))\,\dd V_g<\infty.
 \label{eq:stress-charge-weight}
\end{equation}
The last integrability follows from
$1-t(\zeta)=O(\zeta^{3/2})$ and $s_0>4/3$.
For every smooth function $v$ equal to $-M_e$ near $S$ and zero
outside a compact set, each regular level $v=q$, $-M_e<q<0$, obeys
\begin{equation}
 \int_{v=q}\zeta\,\dd A_g
 \ge(1-o_{R_0}(1))A_*.
 \label{eq:stress-weighted-level-area}
\end{equation}
Here is the topological and geometric comparison for this assertion.
Take the closure $D$ of the component of $\{v>q\}$ containing the
end.  It is a connected closed end-side domain, with boundary a
subset of the level; it misses a neighborhood of all of $S$.
Choose a sphere $r=R_1$, with $R_1\in[R_0/2,R_0]$, transverse to its
boundary, and adjoin to $D$ the whole region $r\ge R_1$.
The union is still connected because the two domains have a common
far-end region.  Each newly exposed sphere-cap point lies on a
radial ray meeting the old boundary at some $r\ge R_1$: initially
the ray is outside $D$, and eventually it lies in $D$.
Radial projection onto $r=R_1$ has upper two-dimensional Jacobian
\[
 (1+o_{R_0}(1))(R_1/r)^2
 \le(1+o_{R_0}(1))\zeta(r).
\]
The first estimate is uniform over tangent two-planes by the
asymptotic tensor bounds; the second uses $s_0<2$ and $R_1\le R_0$.
The area formula bounds the new caps by the weighted old boundary
outside the sphere.  Inside, the weight is one.  The transverse
corners can be rounded toward the side enlarging this union, in
shrinking tubular neighborhoods of their intersection curves; the
added area tends to zero.  The rounded exterior is connected and
its full intrinsic boundary is an admissible cut in the definition
of $A_*$.  This proves \eqref{eq:stress-weighted-level-area} without
any connectedness assertion about the level itself.

Coarea therefore gives
$\int\zeta|dv|\ge M_e(1-o_{R_0}(1))A_*$.
The same holds for $u$.  To justify approximation, first replace
$u$ by its constant trace in a shrinking collar of $S$, interpolating
in a collar of comparable width.  The error in its weighted gradient
integral tends to zero by the bounded gradient and constant trace.
Next multiply by a cutoff tending to zero in $R<r<2R$ and set it
equal to zero beyond $2R$.  The height and gradient decay make both
the lost tail and cutoff error $O(R^{1-s_0})$, tending to zero.
Smooth the remaining compactly supported transition while retaining
the constant boundary neighborhoods.  This gives convergence in
the required weighted $W^{1,1}$ norm, and coarea applies to the
smooth approximants.

Finally, the definition of $L$ gives $L(w)\ge\zeta w+t(\zeta)$,
whereas $L(w)-1\le we(w)$ since $t(e)\le1$.  Combining these facts
with \eqref{eq:stress-energy-charge} yields
\begin{equation}
 F_{M_e}\ge(1-o_{R_0}(1))A_*-C_\zeta/M_e.
 \label{eq:stress-charge-lower-bound}
\end{equation}
Let $M_e\to\infty$ at fixed $R_0$ and then $R_0\to\infty$.
This proves \eqref{eq:stress-charge-saturation} and completes the
proposition.
\end{proof}

\subsection{Smoothing the synthetic tensor and attaching an exterior}

\begin{proposition}[Compact stress data and AF attachment]
\label{prop:af-attachment}
Let $g$ be the fixed approximating metric above.  For every
$\varepsilon_0>0$ there are a sufficiently large coordinate sphere
$\mathcal C$, its compact connected inner region $\Omega_i$, a
smooth symmetric tensor $K$ on $\Omega_i$, and a smooth scalar-flat
\AF\ exterior $(\Omega_o,g_o)$ attached along $\mathcal C$, with
the following properties.  Put $\tau_K=\tr_g K$ and
$k_T=\tr_{\mathcal C}K$.  Then
\begin{align}
 K&<0,\qquad
 R_g+\tau_K^2-|K|_g^2
   -2(\divg_gK-d\tau_K)\cdot v>0
       &&(|v|_g\le1),
 \label{eq:stress-synthetic-DEC}\\
 H_i&>|k_T|,\qquad
 H_o=\sqrt{H_i^2-k_T^2}>0,
 \label{eq:stress-matching-curvatures}\\
 m_o&\le p_0(g)-\tfrac12 a^3+\varepsilon_0.
 \label{eq:stress-attachment-budget}
\end{align}
Here $H_i$ is the inner metric's mean curvature on $\mathcal C$,
and $H_o$ is the outer metric's inner-boundary mean curvature, both
in the direction toward the respective end.  The induced metrics
agree.  The exterior is complete with its boundary included and
has one end, with $g_o-\delta=O_k(|x'|^{-1})$ for every fixed $k$,
where an ordinary derivative improves decay by one power.
Projection from $\Omega_o$ to $\mathcal C$ along its parallel rays
has upper area Jacobian at most one on every tangent two-plane.
\end{proposition}

\begin{proof}
Choose an electrostatic height $M_e$ large enough that the preceding
proposition gives
\begin{equation}
 \langle e_0\rangle\ge a^2-o(1),\qquad
 \langle f\rangle=(4\pi)^{-1}\int_{\Sph^2}f\,\dd A_\sigma,
 \label{eq:stress-charge-average}
\end{equation}
where the error can be arbitrarily small.  On the whole exterior
first define a continuous tensor $K^{(0)}$ by
\begin{equation}
 K^{(0)}-(\tr_gK^{(0)})g=2\Pi,
 \qquad K^{(0)}=2\Pi-(\tr_g\Pi)g.
 \label{eq:stress-synthetic-tensor}
\end{equation}
Its eigenvalues are $t-2l=-t-2we<0$ in the gradient direction and
$-t<0$ transversely.  Directly from
\eqref{eq:stress-algebraic-identity},
\begin{equation}
 (\tr_gK^{(0)})^2-|K^{(0)}|_g^2=6,
 \qquad
 \divg_g\big(K^{(0)}-(\tr_gK^{(0)})g\big)=0
 \label{eq:stress-synthetic-identities}
\end{equation}
in distributions.  For example, if $s_\Pi=\tr_g\Pi$, then
$\tr_gK^{(0)}=-s_\Pi$ and
$|K^{(0)}|^2=4|\Pi|^2-s_\Pi^2$, giving the first identity.

On a coordinate sphere of radius $r$, the induced metric is
$\phi^4r^2\sigma$, hence its Gauss curvature is positive for large
$r$.  The conformal mean-curvature formula gives
\begin{align}
 H_i
 &=\phi^{-2}\left(\frac{2\sqrt{1+r^2}}r
       +4\sqrt{1+r^2}\,\partial_r\log\phi\right)\notag\\
 &=2+r^{-2}-16\alpha r^{-3}+o(r^{-3}).
 \label{eq:stress-inner-H-expansion}
\end{align}
The relation \eqref{eq:stress-synthetic-tensor} gives exactly
$\tr_{\mathcal C}K^{(0)}=-2\Pi(\nu,\nu)$, so
\eqref{eq:stress-end-asymptotics} yields
\begin{equation}
 k_T^{(0)}=-2+e_0^{3/2}r^{-3}+o(r^{-3}).
 \label{eq:stress-tangential-K-expansion}
\end{equation}
All these remainders are uniform in angle.  Consequently
$H_i>|k_T^{(0)}|$ on every sufficiently large sphere, and
\begin{equation}
 \sqrt{H_i^2-(k_T^{(0)})^2}
 =\frac2r+\frac{-16\alpha+e_0^{3/2}}{r^2}+o(r^{-2}).
 \label{eq:stress-outer-H-expansion}
\end{equation}
Indeed the expression under the root is
$4r^{-2}+(-64\alpha+4e_0^{3/2})r^{-3}+o(r^{-3})$.

We explain carefully the smoothing used after fixing such a sphere
$\mathcal C$.  On the compact $\Omega_i$, continuity and strict
positivity give positive lower bounds for the eigenvalues of $\Pi$,
the absolute negative eigenvalues of $K^{(0)}$, and $R_g+6$.
There are smooth tensors $\Pi_h$ on $\Omega_i$, up to its boundary,
such that
\begin{equation}
 \|\Pi_h-\Pi\|_{C^0}\longrightarrow0,
 \qquad \|\divg_g\Pi_h\|_{C^0}\longrightarrow0.
 \label{eq:stress-divergence-smoothing}
\end{equation}
Here is the needed commutator proof, including the boundary.
In an interior chart convolve the tensor components with a smooth
kernel $\eta_h$.  A divergence term has form
$a(x)\partial_jT+b(x)T$, with smooth coefficients.  For the principal
part, integration by parts in the commutator produces
$(a(y)-a(x))\partial\eta_h(x-y)$ and a term with
$(\partial a)(y)\eta_h(x-y)$.  Together these cancel on constant
$T$.  After subtracting $T(x)$ their integral is bounded by a fixed
constant times the modulus of continuity of $T$ on a ball of radius
$O(h)$, and hence tends uniformly to zero.  The lower-order terms
commute up to the same type of error.

In a boundary chart evaluate the convolution instead at
$x+hc_0n_0$, with $n_0$ a constant inward coordinate direction and
$c_0$ large enough that the kernel stays on the original side.
The same proof applies: $|a(y)-a(x)|=O(h)$ throughout its support,
and the combined commutator still vanishes on a constant tensor.
It therefore extends the uniform estimate to the boundary without
extending the original stress as a divergence-free tensor.  Take
these local approximations on slightly larger charts and recombine
with a smooth partition of unity.  The extra differentiated
partition terms tend to zero, because their limiting sum is
$(\sum d\chi)\Pi=0$.  This proves
\eqref{eq:stress-divergence-smoothing}.

Set $K_h=2\Pi_h-(\tr_g\Pi_h)g$.  Then
$\divg K_h-d\tr K_h=2\divg\Pi_h\to0$ uniformly, and the algebraic
part $R_g+(\tr K_h)^2-|K_h|^2$ tends uniformly to $R_g+6>0$.
For sufficiently small $h$, \eqref{eq:stress-synthetic-DEC} follows
uniformly for $|v|\le1$, as does $K_h<0$.  The strict gap
$H_i>|k_T^{(0)}|$ at the now fixed sphere also persists.
Write $K$ for this smooth tensor.  The error in any integral
involving $\sqrt{H_i^2-k_T^2}$ can be made arbitrarily small at this
fixed sphere, because the square root is smooth on its positive
gap.  This order of choices is essential; we require no smoothing
estimate uniform as the sphere tends to infinity.

The positive-Gauss metric on $\mathcal C\simeq\Sph^2$ has a smooth
strictly convex Euclidean isometric embedding by the smooth Weyl
embedding theorem \citep{Nirenberg1953}.  This is precisely the
spherical positive-curvature version also stated in the introduction
of \citet{ShiTam2002}.  Denote its outward Euclidean mean curvature
by $H_E$.  We will construct the Shi--Tam scalar-flat extension with
the smooth positive prescribed boundary curvature
$H_o=\sqrt{H_i^2-k_T^2}$.  The required estimate is
\begin{equation}
 m_o\le\frac1{8\pi}\int_{\mathcal C}(H_E-H_o)\,\dd A.
 \label{eq:stress-shitam-mass-bound}
\end{equation}
Before proving this extension assertion, we evaluate its right side
and record the exact coefficient of the charge term.

Scale the induced metric on $\mathcal C$ by $r^{-2}$.
Its area is $4\pi+O(r^{-3})$ and its Gauss curvature is
$1+O(r^{-3})$, with smooth angular control.  For its convex Euclidean
embedding choose an interior point as origin and use Gauss-map
coordinates $\omega$.  Let $l(\omega)>0$ be the support function.
Then the embedding is $l\omega+\nabla_\sigma l$, its inverse shape
operator is $\nabla_\sigma^2l+l\sigma$, and
\begin{equation}
 \int H_E^{\mathrm{scaled}}\,\dd A
   =\int_{\Sph^2}\tr_\sigma(\nabla_\sigma^2l+l\sigma)\,\dd A_\sigma
   =2\int_{\Sph^2}l\,\dd A_\sigma.
 \label{eq:stress-support-mean}
\end{equation}
The scaled enclosed volume is
$\frac13\int l\,dA$, and $dA=(1+O(r^{-3}))dA_\sigma$ by the Gauss
curvature estimate.  Euclidean isoperimetry bounds this volume by
$4\pi/3+O(r^{-3})$.  Positivity of $l$ now implies
$\int l\,dA_\sigma\le4\pi+O(r^{-3})$.  Scaling back in
\eqref{eq:stress-support-mean} gives
\begin{equation}
 \int_{\mathcal C}H_E\,\dd A\le8\pi r+O(r^{-2}).
 \label{eq:stress-Euclidean-H-bound}
\end{equation}
This uses no quantitative stability theorem for the Weyl embedding.

Before smoothing the stress,
\eqref{eq:stress-outer-H-expansion} and
$dA=\phi^4r^2dA_\sigma$ give
\[
 \int_{\mathcal C}H_o\,\dd A
 =8\pi r+\int_{\Sph^2}(-16\alpha+e_0^{3/2})\,\dd A_\sigma+o(1).
\]
The smoothing error in this equality can be made arbitrarily small
after fixing $r$.  Moreover, direct substitution of
$g-b=(4\alpha r^{-3}+O(r^{-3-\xi}))b$ in the mass flux gives
\begin{equation}
 p_0(g)=8\langle\alpha\rangle.
 \label{eq:stress-conformal-mass}
\end{equation}
For clarity, if $e=fb$, the flux one-form is
$-2(V_0\,df-f\,dV_0)$; for $f=4\alpha r^{-3}$ its radial sphere
integral tends to $32\int\alpha\,dA_\sigma$.
Consequently \eqref{eq:stress-Euclidean-H-bound} gives
\begin{align}
 \frac1{8\pi}\int_{\mathcal C}(H_E-H_o)\,\dd A
 &\le p_0(g)-\tfrac12\langle e_0^{3/2}\rangle+o(1)\notag\\
 &\le p_0(g)-\tfrac12\langle e_0\rangle^{3/2}+o(1)
 \le p_0(g)-\tfrac12a^3+o(1).
 \label{eq:stress-cubic-budget}
\end{align}
The second inequality is Jensen's inequality.  If $a=0$, its charge
term is simply nonnegative.  In all cases first choosing $M_e$,
then $r$, then the compact smoothing tolerance makes the final error
smaller than the prescribed $\varepsilon_0$.

\emph{The scalar-flat extension and its mass.}
We include the extension argument to fix both the existence hypotheses
and the normalization in \eqref{eq:stress-shitam-mass-bound}.  It is
the quasi-spherical construction of \citet[Section~2, Theorem~2.1,
and Lemma~4.2]{ShiTam2002}.  The input here is exactly a smooth
strictly convex Euclidean sphere and a smooth positive prescribed
mean curvature; no compact fill-in is used.

Let $t\ge0$ be Euclidean distance along the outward normals to the
convex embedding.  The parallels foliate its entire exterior.  Write
the Euclidean metric as $dt^2+\gamma_t$, and denote the Euclidean
mean curvature and surface scalar curvature by $H_E(t)$ and $R_t$.
Both are positive.  If $\lambda_1,\lambda_2>0$ are the initial
principal curvatures, then those at time $t$ are
$\lambda_j/(1+t\lambda_j)$, whence
\begin{equation}
 \frac{R_t}{2H_E(t)}
 =\frac1{2t+c(\omega)},\qquad
 c=\lambda_1^{-1}+\lambda_2^{-1}>0.
 \label{eq:stress-parallel-curvature}
\end{equation}
Solve the scalar parabolic equation
\begin{equation}
 2H_E\partial_tU
 =2U^2\Delta_{\gamma_t}U+(U-U^3)R_t,
 \qquad U(0,\cdot)=H_E(0,\cdot)/H_o.
 \label{eq:stress-quasispherical-equation}
\end{equation}
Smooth positive initial data give a local solution by strictly
parabolic theory.  The reaction term is nonpositive above one and
nonnegative below one.  The maximum principle therefore keeps $U$
between fixed positive constants.  On every finite time interval
the equation is uniformly parabolic with smooth background geometry.
Scalar parabolic H\"older estimates for bounded solutions, followed
by Schauder estimates and bootstrapping, give continuation for all
finite times.

For the asymptotic estimate, choose positive constants respectively
above and below the entire initial range and one.  Solve, for each
of these constants as initial value,
\begin{equation}
 \dot z=(z-z^3)/(2t+c_{\max}),\qquad
 c_{\max}=\max c.
 \label{eq:stress-parabolic-barriers}
\end{equation}
By \eqref{eq:stress-parallel-curvature} these are respectively upper
and lower barriers: the sign of $z-z^3$ reverses the inequality in
the upper case.  Since
\[
 z^{-2}-1=(z(0)^{-2}-1)\frac{c_{\max}}{2t+c_{\max}},
\]
comparison gives $U-1=O((1+t)^{-1})$.
Put $s=\log(1+t)$ and $\widehat\gamma_t=(1+t)^{-2}\gamma_t$.
In the Gauss-map parametrization the parallel embedding is
$X(\omega)+t\omega$, so these scaled metrics and all their
derivatives have bounded geometry.  Equation
\eqref{eq:stress-quasispherical-equation} in $s$ is uniformly
parabolic.  Its local H\"older and then Schauder estimates on fixed
$s$-strips first bound all derivatives of $U$ uniformly.  Next apply
the linear estimates to $U-1$ with the actual coefficients of this
solution; its equation is homogeneous, with a bounded smooth
zeroth-order coefficient.  The supremum bound just proved gives
the same $O(e^{-s})$ decay at every fixed derivative order.  This
proves the stated ordinary AF derivative bounds in the Euclidean
coordinates $x'=X(\omega)+t\omega$.

Set
\begin{equation}
 g_o=U^2dt^2+\gamma_t.
 \label{eq:stress-af-extension-metric}
\end{equation}
It has the required induced metric and inner mean curvature
$H_E/U=H_o$.  Its scalar curvature is, by the hypersurface Gauss
and mean-variation formulas,
\begin{equation}
 R_{g_o}=R_t(1-U^{-2})+2H_EU^{-3}\partial_tU
                    -2U^{-1}\Delta_{\gamma_t}U=0,
 \label{eq:stress-af-scalar-zero}
\end{equation}
where the final equality is exactly
\eqref{eq:stress-quasispherical-equation}.  The positive bounds on
$U$ and the Euclidean completeness of the parallel exterior give
completeness with the boundary included.  There is precisely one
end, and $g_o-\delta=(U^2-1)dt^2$ has the stated decay.

Define
\[
 \mathcal M(t)=\int_{\mathcal C_t}H_E(1-U^{-1})\,\dd A_{\gamma_t}.
\]
The Euclidean identities
$\partial_tH_E=-|\mathrm{II}_E|^2$,
$\partial_t dA=H_EdA$, and
$H_E^2-|\mathrm{II}_E|^2=R_t$, together with
\eqref{eq:stress-quasispherical-equation}, give
\begin{align}
 \mathcal M'(t)
 &=\int_{\mathcal C_t}\left[
       R_t(1-U^{-1})+H_EU^{-2}\partial_tU\right]\,\dd A\notag\\
 &=-\frac12\int_{\mathcal C_t}R_t(U+U^{-1}-2)\,\dd A\le0.
 \label{eq:stress-shitam-monotonicity}
\end{align}
Here the integral of $\Delta_{\gamma_t}U$ vanishes.  The asymptotic
bounds show that $\mathcal M$ is bounded and that its derivative
is integrable at infinity, so it has a finite limit.

We identify the limit with the ADM mass in the normalization
\eqref{eq:adm-mass}.  For the purely normal Euclidean metric error
$e'=(U^2-1)dt^2$, the ADM flux through the parallel surface is exactly
\[
 \int_{\mathcal C_t}
   (\divg_\delta e'-d\tr_\delta e')(\partial_t)\,\dd A
 =\int_{\mathcal C_t}H_E(U^2-1)\,\dd A.
\]
The normal derivatives of $U^2-1$ cancel, and
$\nabla_{\partial_t}\partial_t=0$ in the Euclidean metric, leaving
only the tangential divergence $H_E$.  The scalar-curvature
linearization and $R_{g_o}=0$ show that the divergence of this
Euclidean flux is $O(|x'|^{-4})$.  Its integral between a large
parallel and a comparable large round sphere tends to zero; thus
it computes the ordinary ADM mass.  Finally
\[
 (U^2-1)-2(1-U^{-1})=O((U-1)^2)
\]
and $H_E\,\Area(\mathcal C_t)=O(t)$, so
\begin{equation}
 16\pi m_o
 =\lim_{t\to\infty}\int_{\mathcal C_t}H_E(U^2-1)\,\dd A
 =2\lim_{t\to\infty}\mathcal M(t)
 \le2\mathcal M(0).
 \label{eq:stress-adm-identification}
\end{equation}
As $H_E/U(0)=H_o$, this proves
\eqref{eq:stress-shitam-mass-bound}.  Combining it with
\eqref{eq:stress-cubic-budget} gives
\eqref{eq:stress-attachment-budget}.

Lastly, convexity gives
$\gamma_t(V,V)=\gamma_0((\Id+t\mathrm{II}_E^{\sharp})V,
(\Id+t\mathrm{II}_E^{\sharp})V)\ge\gamma_0(V,V)$.
The metric in \eqref{eq:stress-af-extension-metric} has no mixed
normal--tangential terms.  Its projection along the parallels to
$(\mathcal C,\gamma_0)$ therefore has operator norm at most one,
and hence upper two-dimensional Jacobian at most one.  This is
the area comparison required in the later transfer argument.
\end{proof}

All data in Proposition~\ref{prop:af-attachment} are now fixed before
the compact graph parameters are chosen.  On $S$ and $\mathcal C$
we use the base normals directed toward the end; derivatives and
fluxes on the two sides of $\mathcal C$ use $g$ and $g_o$ respectively.
Their boundary area forms agree.  In the next section the notation
$\tau=\tr_gK$ refers to the synthetic trace, not to the original
asymptotic decay exponent.
